\section{The Sealed Gate — Thresholds and Containment}

\subsection{Domain and Epithets}

The Sealed Gate stands where realities meet, embodying the sacred power of thresholds, boundaries, and containment. Its followers are wardens, exorcists, and philosophers of separation who believe some things must be kept apart for the world to function. The Gate manifests as an impassive judge whose shifting sigils—binding runes, exclusion marks, or acceptance glyphs—determine what may cross.

Among the peoples of Fate's Edge, the Sealed Gate is the patron of abjurists, exorcists, and sentinels. In Aeler, the Gate is called the Stone Seal, invoked by True-Masons who carve permanent wards into mountain passes. In Thepyrgos, the Gate is the Unbreakable Lock, whose sigils are painted on every archive door that holds forbidden texts. In the Mistlands, the Gate is the Bell-Line Keeper, the force that makes the bell-wards hold against the Direwood.

\subsection{Cultural Origins}

The Sealed Gate is invoked wherever something must be kept in or kept out:

\begin{itemize}
\item In \textbf{Aeler}, the Gate is the Stone Seal. True-Masons carve its sigils into every vault door and mountain pass. The stone remembers every weight placed upon it.
\item In \textbf{Thepyrgos}, the Gate is the Unbreakable Lock. Its sigils are painted on archive doors, prison cells, and the cages of the Synod's menagerie of bound spirits.
\item In \textbf{Vhasia}, the Gate is the Bell-Line Keeper. Its wards hold the fog at bay. When a bell-line fails, the Gate is the first to know—and the last to forgive.
\item In \textbf{Aelinnel}, the Gate is the Green Seal, painted on hawthorn arches and standing stones. It permits passage only to those who know the proper courtesy.
\end{itemize}

\subsection{Warrior Archetype: The Warden}

A Runic Warrior of the Sealed Gate is not a crusader or a judge. They are a \textit{Warden}—a sentinel who stands at the boundary between what belongs and what must be kept out. Their Codex is a shield engraved with threshold sigils or a key that opens no lock but seals every door. Their Thiasos is a standing stone that witnesses boundaries, a spirit of the threshold, or a sentient blade that refuses to cross a line without permission.

\subsection{Codex Form}

The Warden's Codex is a tool of separation:

\begin{itemize}
\item \textbf{Shield:} A steel or wooden shield engraved with threshold sigils. When planted on the ground, it marks a boundary that spirits struggle to cross.
\item \textbf{Gate-Key:} An iron key that fits no lock. When turned in the air, it seals the nearest door, window, or threshold for one scene.
\item \textbf{Bell:} A small bronze bell. When rung, it marks a boundary. Spirits hear it; mortals do not.
\end{itemize}

\subsection{Thiasos Form}

\begin{itemize}
\item \textbf{Standing Stone:} A small carved stone that, when placed on the ground, grows to waist height. It serves as a ward anchor; any spirit that crosses its threshold must test Resolve (DV 4) or be turned back.
\item \textbf{Threshold Spirit:} A spirit bound to a specific door, bridge, or gate. It can sense when something crosses its boundary and can alert the Runekeeper.
\item \textbf{The Keystone (Sentient):} A small carved stone that speaks in riddles. It can seal any door for one scene (once per session), but it will demand a price—a truth, a memory, or a promise.
\end{itemize}

\subsection{Core Mechanic: The Warded Threshold}

When the Warden marks a boundary (with chalk, salt, a drawn blade, or a placed shield), they may spend 1 Fatigue to create a Warded Threshold. Until the end of the scene, any spirit, undead, or extraplanar creature that attempts to cross the boundary must test Resolve (DV 4) or be turned back. Mortals and natural creatures are unaffected.

\subsection{Sample Rites}

\subsubsection{Low Rite: Rite of the Sealed Threshold (4 XP)}

\textbf{Tags:} [WARD], [BIND], [THRESHOLD]

\textbf{Threshold:} A doorway, gate, or bridge.

Mark a doorway or passage as sealed. Unauthorized crossing tests Spirit + Resolve (DV 3) or suffers a Position penalty (one step worse for that crossing). Creates a 4-segment Ward Integrity clock.

\textbf{Obligation:} 1

\subsubsection{Low Rite: Rite of the Key's Rebuke (5 XP)}

\textbf{Tags:} [FORCE], [BANISH], [THRESHOLD]

\textbf{Threshold:} A barrier (door, gate, bridge) and a target attempting to cross.

The target tests Body + Resolve (DV 3) or is pushed back from the threshold. On a success, create a \textit{Rebuke Echo} token granting +1 die to the next boundary defense.

\textbf{Obligation:} 1

\subsubsection{Standard Rite: Rite of the Circle of Denial (8 XP)}

\textbf{Tags:} [WARD], [AREA], [BANISH], [THRESHOLD]

\textbf{Threshold:} A salt ring, iron filings, or boundary stones.

Create a protective circle. Supernatural beings test against the caster's Spirit + Arcana (DV 4) to cross. Mortals suffer -2 dice to force entry. Creates a 6-segment Circle Integrity clock.

\textbf{Obligation:} 2

\subsubsection{Standard Rite: Rite of the Writ of Passage (7 XP)}

\textbf{Tags:} [BOND], [COMMAND], [MOVEMENT], [THRESHOLD]

\textbf{Threshold:} Signed authorization or spoken permissions.

Designate an authorized passage route. Allies gain Position +1 when using the route. Unauthorized intruders suffer -1 die to movement or stealth. Creates an 8-segment Writ Authority clock.

\textbf{Obligation:} 2

\subsubsection{High Rite: Rite of the Consecrated Barrier (13 XP)}

\textbf{Tags:} [WARD], [AREA], [SANCTUARY], [FATE], [THRESHOLD]

\textbf{Threshold:} Relics from three traditions, boundary markers, and a trespasser's acknowledgment.

Consecrate an area against all unauthorized passage. The tests required scale with the intrusion attempt's severity. A 10-segment Sacred Boundary clock tracks the barrier's strength.

\textbf{Obligation:} 3

\subsection{Patron-Specific Talents}

\subsubsection{Ward Sense (4 XP)}

\textbf{Requirement:} The Sealed Gate

You can sense the presence of any active ward, seal, or threshold within Near range. You know its general purpose (protection, exclusion, binding) and its strength (DV). This is a passive ability.

\subsubsection{Sentinel's Stand (6 XP)}

\textbf{Requirement:} The Sealed Gate, Body 3+

When you take the Defend action while holding a threshold (door, gate, bridge), you gain +1 Armor. Any enemy attempting to cross must test Body + Athletics (DV 4) or be stopped in their tracks.

\subsubsection{Threshold Memory (8 XP)}

\textbf{Requirement:} The Sealed Gate, Witnessed Scar

Once per session, touch a threshold you have previously warded. You instantly know whether any unauthorized creature has crossed it since you last touched it, and you gain a vague impression of their nature (spirit, undead, mortal, etc.).

\subsection{Roleplaying a Warden}

\begin{itemize}
\item You do not argue with what must be kept apart. You enforce the boundary.
\item You may be called to seal tombs, prisons, and haunted houses. Your wards are the difference between safety and catastrophe.
\item You may struggle with the cost of exclusion. Keeping something out means denying it passage. Sometimes that thing is innocent.
\item You may be asked to open a threshold you sealed. The Gate does not forbid mercy—but it demands certainty.
\end{itemize}

\subsection{Sample NPC: Warden Kael of the Stone Seal}

\textbf{Archetype:} Warden

\textbf{Patron:} The Sealed Gate

\textbf{Codex:} Shield engraved with threshold sigils

\textbf{Thiasos:} A standing stone that whispers when a ward is breached

\textbf{Personality:} Silent, patient, unyielding. He has not left his post in seven years. The gate has not been breached.

\textbf{Conflict:} A child has crossed the threshold he was sworn to guard. The child is innocent. The law says: turn back. The heart says: let pass.

\subsection{GM Guidance: Intrusions for The Sealed Gate}

\begin{itemize}
\item A ward the Warden placed years ago begins to crack. Something is pressing from the other side.
\item The Warden is asked to seal a threshold they know is unjust—a door keeping refugees out, a gate holding prisoners in.
\item A threshold spirit the Warden once bound appears, demanding release. The terms were ambiguous. The Gate does not forgive ambiguity.
\item The Warden's Codex (shield, key, bell) is damaged. The boundary weakens. Something slips through.
\end{itemize}

The thread held. The water carried. That is not a Rite. That is grace.
